\input{lib/dissertationelse.tex}
\input{lib/dissertationexclude.tex}
\input{lib/dissertationonly.tex}

\section{Lineage History}

\input{fig/keyfig.tex}

In brief, specimens from early stints largely grew as unicellular or small multicellular groups (morphs $a$, $b$).
Then, the specimen from stint 14 grew as larger, symmetrical groups (morph $d$).
At stint 15, a distinct, asymmetrical horizontal bar morphology evolved (morph $e$).
%todo structure above is confusing
% Consistently left/right, indicating that somehow broke symmetrical of the simulation.
At stint 45, a delayed secondary spurt of group growth in the vertical direction arose (morph $g$).
This morphology was sampled frequently until stint 60, when morph $e$ began to be sampled primarily again.
However, morph $g$ was observed as late as stint 90.


Across the case study lineage, we describe an evolutionary sequence of ten qualitatively distinct multicellular morphologies (Table \ref{tab:morph_descriptions}).
We performed a qualitative survey of the evolved life histories along the evolutionary timeline by reviewing video animations of the representative specimen from each stint in monoculture.

\input{table/morph_descriptions.tex}


Table \ref{tab:morph_descriptions} summarizes the ten morphological categories we grouped specimens into.


\subsection{Notable Events}

\begin{itemize}
\item \textbf{Stint 2}:
\item \textbf{Stint 3}:
\item \textbf{Stint 14}:
\item \textbf{Stint 15}:
\item \textbf{Stints 16 through 39}:
\item \textbf{Stints 27 and 29}:
\item \textbf{Stint 45}:
\item \textbf{Stint 50}:
\item \textbf{Stints 74, 75, 77, and 92}:
\item \textbf{Stint 93}:
\item \textbf{Stints 95 and 96}:
\end{itemize}

\subsection{Evolutionary History}

% https://mybinder.org/v2/gh/mmore500/dishtiny/e17e6d5e258b7aacac72d44922008ab14e80e182?filepath=binder%2Fbucket%3Dprq49%2Fa%3Dall_stints_all_series_profiles%2Bendeavor%3D16%2Fcase_study_16005.ipynb
Due to the parallel, asynchronous nature of the experimental framework, we did not perform perfect phylogeny tracking \citep{moreno2024analysis}.
\dissertationonly{Chapter \ref{ch:distributed-phylogeny} discusses challenges parallelizing perfect phylogeny tracking in depth.}

Instead, we opted to track the total number of ancestors seeded into stint 0 with extant descendants.
At the end of stints 0 and 1, three distinct original phylogenetic roots were present in the population.
From stint 2 onward, only two distinct original phylogenetic roots were present.

% https://mybinder.org/v2/gh/mmore500/dishtiny/e17e6d5e258b7aacac72d44922008ab14e80e182?filepath=binder%2Fbucket%3Dprq49%2Fa%3Dall_stints_all_series_profiles%2Bendeavor%3D16%2Fcase_study_16005.ipynb
We performed follow-up analyses on specimens sampled from the lowest original phylogenetic root ID present in the population.%
\footnote{
This approach was designed to choose an arbitrary strain as focal.
Barring extinction, that same strain will then be identified as focal consistently across subsequent stints.
Phylogenetic root ID had no functional consequences; it is simply an arbitrary basis for focal strain selection.
}
For the first two stints, the focal strain was root ID 2,378.
During stint 2, original phylogenetic root 2,378 went extinct.
So, all further follow-up analyses were sampled from descendants of ancestor 12,634.

We also tracked the number of genomes reconstituted at the outset of each stint that had extant descendants at the end of that stint.
This count grows from approximately 10 around stint 15 to upwards of 30 around stint 40 (Supplementary Figure \ref{fig:phylogeny:stint_roots}).
Among descendants of the lowest original phylogenetic root, the number of independent lineages spanning a stint also increases from around 5 to around 15
(Supplementary Figure \ref{fig:phylogeny:lowestroot_stint_roots}).
This decrease in phylogenetic consolidation on a stint-by-stint basis correlates with the waning number of simulation updates performed per stint (Supplementary Figures \ref{fig:phylogeny:updates_vs_stint_roots} and \ref{fig:phylogeny:log_updates_vs_stint_roots}).
More complete phylogenetic data will be necessary in future experiments to address questions about the possibility of long-term stable coexistence beyond the two strains supported under the explicit diversity maintenance scheme.

% https://mybinder.org/v2/gh/mmore500/dishtiny/e17e6d5e258b7aacac72d44922008ab14e80e182?filepath=binder%2Fbucket%3Dprq49%2Fa%3Dall_stints_all_series_profiles%2Bendeavor%3D16%2Fcase_study_16005.ipynb
On the specimen from stint 100 used in the final case study, an evolutionary history of 20,212 cell generations had elapsed.
Of these cellular reproductions, 11,713 (58\%) had full kin group commonality, 7,174 had partial kin group commonality (35\%), and 1,325 had no kin group commonality (7\%).
On this specimen, 1,672 mutation events had elapsed.
During these events, 7,240 insertion-deletion alterations had occurred and 26,153 point mutations had occurred.
This strain experienced a selection pressure of 18\% over its evolutionary history, meaning that only 82\% of the mutations that would be expected given the number of cellular reproductions that had elapsed were present.

% https://github.com/mmore500/dishtiny_event_tag_phylogenetics/commit/ecc2ab8a580b125bea9f661e49cd4d7f34bfa1bb
In order to characterize the evolutionary history of the experiment in greater detail, we performed a parsimony-based phylogenetic reconstruction on the sampled representative specimens from each stint, shown in Supplementary Figure \ref{fig:phylo_parsimony_tree}.
We used genomes' fixed-length blocks of 35 64-bit tags that mediate environmental interactions as the basis for this reconstruction.
These tag blocks underwent bitwise mutation over the course of the experiment.%
\footnote{
In future experiments, we plan to incorporate new methodology for ``hereditary stratigraph'' genome annotations expressly designed to facilitate phylogenetic reconstruction \dissertationelse{(Chapter \ref{ch:distributed-phylogeny})}{\citep{moreno2022hereditary}}.
}
Supplementary Figure \ref{fig:phylo_distance_matrix_heatmap} shows Hamming distance between all pairs of tag blocks.
We additionally tried several other tree inference methods, discussed in supplementary material; however, these yielded lower-quality reconstructions.

Although the phylogeny of stint representatives includes many instances that do not constitute a strict sequential lineage (i.e., each stint's representative descending directly from the preceding stint's representative), we did not observe evidence of long-term coexistence of clades over more than ten stints.

\subsection{Qualitative Morphological Categorizations}


Phylogenetic analysis (Supplementary Figure \ref{fig:phylo_parsimony_tree}) indicates that observations of morph $e$ at stint 53 and onward are instances of secondary loss rather than retention of trait $e$ by a separate lineage coexisting with the lineage expressing morph $g$.
Three separate reversion events from morph $g$ to morph $e$ appear likely.
Interestingly, morph $g$ individuals at stints 89 and 90 appear to represent subsequent trait re-gain after reversion from morph $g$ to morph $e$.

Table \ref{tab:morph_descriptions} provides more detailed descriptions of each qualitative morph category, as well as a video and still image example of each.

Supplementary Table \ref{tab:morph_by_stint} provides morph categorization for each stint as well as links to view the stint's specimen in a video or in-browser web simulation.

% \subsection{Multicellular Phenotypic Traits}

% In order for a transition in individuality, cells have to diminish their own immediate reproductive interests in order to prioritize the interests of the collective.
% In previous work with an earlier version of the system, we've seen multicellular traits evolve such as (blah from abstract) %todo
% \citep{moreno2021exploring}.

% The formation of kin group morphologies in this system does not necessarily imply a transition in individuality.
% It is necessary to test to see if cells within groups are meaningfully cooperating.

% Because cell reproduction is necessarily adversarial in this system, we can test to see whether cells preferentially target or spare their group members.
% A conflict ratio of less than 1 implies that cells preferentially spare their members.
% We observed a  0.5 most of the time.

% Supplementary Figure \ref{fig:conflict:exactly_one} shows the  . %todo
% Supplementary Figure \ref{fig:conflict:at_least_one}
% Supplementary Figure \ref{fig:conflict:exactly_two}.
% However, could be an artifact of cell-offspring cooperation.

% Interestingly, this strain appears to be using ``Is Child Cell Of,'' ``Is Parent Cell Of,'' ``Is Child Group Of,'' and the raw kin group ID of neighbors (but not the raw kin group ID of self).

% %todo change wording, not sure what you're hoping to say with first phrase
% To more conclusively, in future work we plan to look at the reproductive output at the interior of multicellular groups to see if these cells diminish their own reproductive output and to control for cellular relatedness when calculating this element of group cooperation.

% We analyzed other traits characteristic of multicellularity, as well.

% We did not see widespread resource sharing over the evolutionary history.
% Only tiny fractions of cells, less than 1\%, appeared to share tiny amounts of the resource, less than 1\% of the resource required to reproduce per update (Supplementary Figures \ref{fig:resource_sharing:fraction_sharing_monoculture} and \ref{fig:resource_sharing:sharing_amount_monoculture}).
% However, in ten strains sampled along the evolutionary lineage, context-dependent resource-sending behavior contributed to fitness (Supplementary Figure \ref{fig:writable_perturbation}).
% Interestingly, another strain present in the population that was not the subject of analysis appeared to evolve widespread, higher-throughput resource sharing after stint 60 (Supplementary Figures \ref{fig:resource_sharing:fraction_sharing_evolve} and \ref{fig:resource_sharing:sharing_amount_evolve}).

% We observed widespread apoptosis over the evolutionary history.
% As shown in Supplementary Figure \ref{fig:apoptosis:apoptosis_stint}, in most monocultured representative strains, between 10 and 30\% of cell deaths were due to apoptosis.
% However, it is not clear if this apoptosis returned a fitness benefit.
% No context-dependent apoptosis behavior contributed to fitness for any sampled strain over the evolutionary history (Supplementary Figure \ref{fig:writable_perturbation}).
% It is unclear whether context-independent apoptosis might have a fitness benefit.

% The behavior seems like cells from within a group don't explicitly recognize one another according to group membership. However, they do recognize their cell parent/child and arrange themselves in a somewhat linear fashion that minimizes the potential for conflict between cells.
% Overall, there are elements of cooperation but is unclear to what extent this particular strain constitutes a proper fraternal transition in individuality